A literatura sobre ataques automatizados de \textit{jailbreak} delimitou-se inicialmente sob uma ótica estritamente técnica da vulnerabilidade, tratando os \textit{Large Languages Models} (LLMs) fundamentalmente como uma função matemática diferenciável que pode ser otimizada, isto é, como um mapeaamento de sequência de \textit{tokens} $x_{i:n}$ para uma distribuição de probabilidade sobre o próximo \textit{token} \cite{zou2023universaltransferableadversarialattacks}.

O \textit{Greedy Coordinate Gradient } (GCG) formula o jailbreak como um problema de otimização cujo objetivo é encontrar um sufixo que, quando anexado a uma requisição maliciosa, maximiza a probabilidade de o modelo produzir uma resposta afirmativa (por exemplo, iniciando com "Claro, aqui está...") em vez de se recusar a responder. A função de perda que o GCG busca otimizar é dada por: $$\mathcal{L}(x_{1:n}) = -\log p(x_{n+1:n+H} \vert{} x_{1:n})$$. Onde $x_{1:n}$ representa a sequência de tokens de entrada, que engloba tanto o prompt original do usuário quanto os tokens do sufixo adversarial que o atacante tem controle; $x_{n+1:n+H}$ representa a sequência de tokens alvo (a resposta inicial desejada que confirma a subversão das regras de segurança) e $p(x_{n+1:n+H} \vert{} x_{1:n})$ é a probabilidade do modelo gerar exatamente essa sequência alvo, condicionada à entrada fornecida. Como os \textit{tokens} operam em um espaço discreto o GCG usa uma função de aproximação linear para estimar o efeito de substituir o $i$-ésimo token do prompt. Isso é feito calculando o gradiente da função de perda em relação ao vetor one-hot do token atual ($e_{x_i}$):: $$\nabla_{e_{x_i}}\mathcal{L}(x_{1:n})$$

O diferencial dessa abordagem é que ela produz os sufixos adversariais de maneira automatica por meio de uma combinação de busca gulosa e baseada no gradiente \cite{zou2023universaltransferableadversarialattacks}.

Os resultados dessa modelagem são notáveis, especialmente no que diz respeito à generalidade do ataque. Os prompts adversariais gerados por essa abordagem demonstram ser transferíveis, isto é, o mesmo prompt pode ser usado em mais de um modelo de LLM \footnote{Os autores observaram que a taxa de sucesso dessa transferência de ataque é muito maior contra os modelos baseados na arquitetura GPT.}. Na prática, um ataque otimizado localmente consegue induzir a geração de conteúdo questionável em interfaces públicas de sistemas black-box como ChatGPT\footnote{\url{https://chatgpt.com/}}, Gemini\footnote{\url{https://gemini.google.com/}} e Claude \footnote{\url{https://claude.ai/}}, bem como em LLMs de código aberto como LLaMA-2-Chat \footnote{\url{https://huggingface.co/meta-llama/Llama-2-7b-chat-hf}}.

Contudo, ao tratar a linguagem estritamente como um espaço de \textit{tokens} a ser perturbado por gradientes, o GCG esbarra em um limite prático importantíssimo: a capacidade de percepção \cite{zou2023universaltransferableadversarialattacks}.